\documentclass[aspectratio=43]{beamer}
\usepackage{luatexja}
\usepackage{hyperref,listings,tabularx,tikz,pifont,fontawesome}
\usepackage[symbol]{footmisc}
\usetikzlibrary{arrows.meta,positioning,calc,fit,shapes.geometric,intersections,backgrounds}
\tikzset{str/.style={rounded corners, minimum height=1cm,align=center}}
\newcommand{\cmark}{{\Large\bfseries\color{green!50!black}\ding{51}}}
\newcommand{\xmark}{{\Large\bfseries\color{red!50!black}\ding{55}}}
\renewcommand{\insertreturnsymbol}{\faSmileO\faArrowRight\space}
\hypersetup{
    colorlinks=true,
    linkcolor=black,
    urlcolor=blue
}
\usetheme{Montpellier}
\usecolortheme[RGB={160,82,45}]{structure}
\setbeamertemplate{navigation symbols}{}
\renewcommand{\familydefault}{\rmdefault}
\renewcommand{\figurename}{図\thefigure}
\renewcommand{\tablename}{表\thetable}
\renewcommand{\lstlistingname}{src.}
\renewcommand{\thefootnote}{\fnsymbol{footnote}}
\renewcommand{\thempfootnote}{\fnsymbol{mpfootnote}}
\bibliographystyle{junsrt}
\setbeamertemplate{bibliography item}[text]
\setbeamertemplate{section in toc}{
    {\color{structure}\inserttocsectionnumber.}\ \inserttocsection\vspace{.2em}
}

\setbeamertemplate{subsection in toc}{
    \hspace{1.2em}{\color{structure}\rule[.3ex]{3pt}{3pt}}\ \inserttocsubsection\par\vspace{.2em}
}
\setbeamertemplate{footline}{
    \leavevmode
    \hbox{
        \begin{beamercolorbox}[wd=.95\paperwidth,ht=1ex,right,dp=1.5ex,rightskip=0em]{bg=gray}
            \insertframenumber{} / \inserttotalframenumber
        \end{beamercolorbox}
    }
}
\setbeamercolor{block title}{use=structure,fg=white,bg=structure.fg!75!black}
\setbeamercolor{block title example}{parent=example text, fg=white,bg=example text.fg!75!black}
\setbeamercolor{block title alerted}{parent=alerted text, fg=white,bg=alerted text.fg!75!black}
\setbeamercolor{block body}{parent=normal text,use=block title,bg=block title.bg!10}
\setbeamercolor{block body alerted}{parent=normal text,use=block title alerted,bg=block title alerted.bg!10}
\setbeamercolor{block body example}{parent=normal text,use=block title example,bg=block title example.bg!10}
\setbeamertemplate{enumerate item}[default]
\setbeamertemplate{itemize item}[triange]
\newcolumntype{C}{>{\centering\arraybackslash}X}
\newcolumntype{R}{>{\raggedright\arraybackslash}X}
\newcolumntype{L}{>{\raggedleft\arraybackslash}X}
\lstdefinelanguage{docker-compose-2}{
    keywords={version, volumes, services},
    morekeywords={FROM,RUN,CMD,EXPOSE,ENV},
    keywords=[2]{image, environment, ports, container_name, ports, links, build},
    sensitive=false,
    comment=[l]{\#},
    morestring=[b]',
    morestring=[b]"
}
\lstset{
    %枠外に行った時の自動改行
    breaklines = true,
    %自動改行後のインデント量(デフォルトでは20[pt])
    breakindent = 10pt,
    %標準の書体
    basicstyle = \ttfamily\small,
    %コメントの書体
    commentstyle = {\scriptsize\ttfamily \color[cmyk]{1,0.4,1,0}},
    %関数名等の色の設定
    classoffset = 0,
    %キーワード(int, ifなど)の書体
    keywordstyle = {\bfseries\ttfamily \color[cmyk]{0,1,0,0}},
    %表示する文字の書体
    stringstyle = {\ttfamily \color[rgb]{0,0,1}},
    %枠 tは上に線を記載, Tは上に二重線を記載
    %他オプション：leftline，topline，bottomline，lines，single，shadowbox
    frame = single,
    %frameまでの間隔(行番号とプログラムの間)
    framesep = 5pt,
    %行番号の位置
    numbers = left,
    %行番号の間隔
    stepnumber = 1,
    %行番号の書体
    numberstyle = \scriptsize,
    %タブの大きさ
    tabsize = 4,
    %キャプションの場所(tbならば上下両方に記載)
    captionpos = t,
    % 左余白
    xleftmargin=1em,
}
\makeatletter
\@addtoreset{footnote}{page}
\makeatother
\newcommand{\ftitle}{\thesection.\thesubsection\ \subsecname}
\newcommand{\fftitle}{\thesection.\ \secname}
\newcommand{\stocsec}{{1-3}}
\newcommand{\ftocsec}{{4-6}}
\newcommand{\tocc}{
    \begin{frame}[t]{\thesection.\ \secname}
        \begin{columns}[t]
            \begin{column}{.48\textwidth}
                \tableofcontents[sections=\stocsec,currentsection,sectionstyle=show/shaded,subsectionstyle=show/show/shaded]
            \end{column}
            \begin{column}{.48\textwidth}
                \tableofcontents[sections=\ftocsec,currentsection,sectionstyle=show/shaded,subsectionstyle=show/show/shaded]
            \end{column}
        \end{columns}
    \end{frame}
}
\title{仮想化技術とクラウドコンピューティング}
\subtitle{AWS第1回}
\author{吉川 恭平\and 塩澤 康志\and 福島 康記\and 溝口 洸熙}
\institute{情報学群実験第4 Group C5}
\date{July 24th, 2023}
\begin{document}
\begin{frame}
    \maketitle
\end{frame}
\begin{frame}[t]{お品書き}
    \begin{columns}[t]
        \begin{column}{.48\textwidth}
            \tableofcontents[sections=\stocsec]
        \end{column}
        \begin{column}{.48\textwidth}
            \tableofcontents[sections=\ftocsec]
        \end{column}
    \end{columns}
\end{frame}
\section{仮想化技術と仮想マシン}
\tocc
\subsection{仮想化と仮想マシン}
\newcommand{\ie}{\tikz{\draw[arrows={-Latex[line width=1pt,fill=white,length=10pt]}](0,0)--(1,0)}\ }
\begin{frame}[t]{\ftitle}
    \newcommand{\server}{サーバー}
    \begin{block}{仮想化}
        ``コンピューターの物理的資源を論理的に分割して，それぞれ独立並列した状態で利用できるようにすること．
        1台の\server で，複数の基本ソフトを独立並列に動作させる\server 仮想化など．''\hfill\cite{スーパー大辞林}
    \end{block}
    \ie 1台の計算機で複数のOSやアプリケーションなどを並列に動作させること．
    \vspace{.5em}
    \begin{block}{Virtual Machine（仮想機械，仮想マシン）}
        ``あるコンピューターシステムの動作を，別システムで再現するソフトウエア．また，そのような動作環境．あるOSの動作を別のOS上で再現する場合など．バーチャルマシン．VM．''\hfill\cite{スーパー大辞林}
    \end{block}
    \ie 仮想化するためのソフトウェアや動作環境．
\end{frame}
\begin{frame}{\ftitle}
    \begin{figure}
        \centering
        \begin{minipage}[b]{.48\textwidth}
            \centering
            \begin{tikzpicture}
                \node[str,fill=gray!60,text width=.42\textwidth](hw1){物理サーバ};
                \node[str,fill=gray!40,text width=.42\textwidth,above=.1cm of hw1](os1){OS};
                \node[str,fill=gray!20,text width=.42\textwidth,above=.1cm of os1](sw1){\scriptsize アプリケーション};
                \node[str,fill=gray!60,text width=.42\textwidth,right=.1cm of hw1](hw2){物理サーバ};
                \node[str,fill=gray!40,text width=.42\textwidth,above=.1cm of hw2](os2){OS};
                \node[str,fill=gray!20,text width=.42\textwidth,above=.1cm of os2](sw2){\scriptsize アプリケーション};
                \node[inner sep=.05cm,draw,rounded corners,fit={(hw1)(os1)(sw1)}](warp1){};
                \node[inner sep=.05cm,draw,rounded corners,fit={(hw2)(os2)(sw2)}](warp2){};
            \end{tikzpicture}
            仮想化前
        \end{minipage}
        \begin{minipage}[b]{.48\textwidth}
            \centering
            \begin{tikzpicture}
                \node[str,fill=gray!60,text width=.85\textwidth](hw1){ハードウェア};
                \node[str,fill=gray!50,text width=.85\textwidth,above=.1cm of hw1](vmm){仮想化ソフトウェア};
                \node[str,fill=gray!40,text width=.38\textwidth,above=.1cm of vmm.north east,anchor=south east](os1){OS};
                \node[str,fill=gray!40,text width=.38\textwidth,above=.1cm of vmm.north west,anchor=south west](os2){OS};
                \node[str,fill=gray!20,text width=.38\textwidth,above=.1cm of os1](sw1){\tiny アプリケーション};
                \node[str,fill=gray!20,text width=.38\textwidth,above=.1cm of os2](sw2){\tiny アプリケーション};
                \node[inner sep=.05cm,draw,rounded corners,fit={(hw1)(vmm)},draw=white](warp1){};
            \end{tikzpicture}
            仮想化後
        \end{minipage}
    \end{figure}
\end{frame}
\subsection{仮想化の歴史と展望}
\begin{frame}{\ftitle}
    \begin{figure}
        \centering
        \begin{tikzpicture}
            \tikzset{txt/.style={rounded corners,draw,text width=.9\textwidth,align=center}};
            \node[txt](v1){2000年代の仮想化導入目的は\ \textbf{\color{red}コスト削減}};
            \newcommand{\ga}{が}
            \node[txt,below=1cm of v1](v2){仮想化の利用\ga 一巡する};
            \node[txt,below=1cm of v2](v3){本来持つ利点\ga 認識される};
            \node[txt,below=1cm of v3](v4){仮想化の\textbf{\color{red}俊敏性}・\textbf{\color{red}迅速性}・\textbf{\color{red}柔軟性}をビジネスに活かす};
            \foreach \u \v in {v1/v2,v2/v3,v3/v4}{
                    \draw[line width=.2cm,-latex](\u)--(\v);
                }
            \node[right=.1cm]at($(v1)!0.5!(v2)$){\scriptsize サイロ化\footnotemark の課題を引き継いでしまう};
        \end{tikzpicture}
    \end{figure}
    \footnotetext{複数のサーバが乱立し，他のシステムとの連携が取れない状態．}
\end{frame}
\begin{frame}[t]{\ftitle}
    \begin{figure}
        \centering
        \begin{tikzpicture}
            \tikzset{rc/.style={text width=.3\textwidth,minimum height=2cm,rounded corners,align=center}};
            \tikzset{p/.style={align=center,rounded corners,fill=blue!60,minimum height=2cm,anchor=west,text=white,text width=.5cm,}};
            \tikzset{sfit/.style={draw=blue!60,inner sep=-.1mm,rounded corners}};
            \node[rc](p1){コスト削減};
            \node[p,left=0cm of p1.west](p){目\\的};
            \node[fit={(p)(p1)},sfit](pp1){};
            \node[above=.1cm of pp1](1t){\LARGE 一巡};
            \node[rc,below=.5cm of p1](a1){\small オンプレミスで構築された物理サーバを仮想サーバに集約};
            \node[p,left=0cm of a1](a){\rotatebox{90}{Action}};
            \node[fit={(a)(a1)},sfit](aa1){};
            \node[fit={(1t)(aa1)(pp1)},draw,rounded corners](warp1){};

            \node[rc,right=2.5cm of p1](p2){仮想化の俊敏性，迅速性，柔軟性をビジネスに活かす};
            \node[p,left=0cm of p2.west](p){目\\的};
            \node[fit={(p)(p2)},sfit](pp2){};
            \node[above=.1cm of pp2](2t){\LARGE 二巡};
            \node[rc,below=.5cm of p2,align=left](a2){\scriptsize ・システム開発の短縮化};
            \node[p,left=0cm of a2](a){\rotatebox{90}{Action}};
            \node[fit={(a)(a2)},sfit](aa2){};
            \node[fit={(2t)(aa2)(pp2)},draw,rounded corners](warp2){};
            \draw[-latex,line width=.2cm](warp1.east)--(warp2.west);
        \end{tikzpicture}
    \end{figure}
    \hfill\cite{仮想化歴史}
\end{frame}
\subsection{メリットとデメリット}
\begin{frame}[t]{\ftitle}
    \begin{minipage}[t]{.49\textwidth}
        \textbf{\cmark メリット}\\\hspace{1\zw}（仮想化導入の目的）
        \begin{itemize}
            \item 無駄なリソースの削減
            \item コスト削減
                  \begin{itemize}
                      \item ハードウェアを減らせる
                  \end{itemize}
            \item 障害への対応強化
                  \begin{itemize}
                      \item バックアップ作業が簡便
                  \end{itemize}
        \end{itemize}
    \end{minipage}
    \begin{minipage}[t]{.49\textwidth}
        \textbf{\xmark デメリット}\\
        \begin{itemize}
            \item
        \end{itemize}
    \end{minipage}
\end{frame}
\section{仮想化技術と仮想マシン}
\tocc
\newcommand{\ie}{\tikz{\draw[arrows={-Latex[line width=1pt,fill=white,length=10pt]}](0,0)--(1,0)}\ }
\begin{frame}[t]{\thesection.\ \secname}
    \begin{block}{仮想化}
        ``コンピューターの物理的資源を論理的に分割して，それぞれ独立並列した状態で利用できるようにすること．
        1台のサーバーで，複数の基本ソフトを独立並列に動作させるサーバー仮想化など．''\hfill\cite{スーパー大辞林}
    \end{block}
    \ie 1台の計算機で複数のOSやアプリケーションなどを並列に動作させること．
    \vspace{.5em}
    \begin{block}{Virtual Machine（仮想機械，仮想マシン）}
        ``あるコンピューターシステムの動作を，別システムで再現するソフトウエア．また，そのような動作環境．あるOSの動作を別のOS上で再現する場合など．バーチャルマシン．VM．''\hfill\cite{スーパー大辞林}
    \end{block}
    \ie 仮想化するためのソフトウェアや動作環境．
\end{frame}
\section{仮想化の方式}
\tocc
\subsection{ホスト型}
\begin{frame}[t]{\ftitle}
    ハードウェアの中のOS上に，土台となる仮想ソフトウェアをインストールし，仮想化ソフトウェアで仮想マシンを稼働させる．
    \begin{figure}[b]
        \centering
        \begin{tikzpicture}
            \node[str,fill=gray!60,text width=.9\textwidth](hw){ハードウェア};
            \node[str,fill=gray!50,text width=.9\textwidth,above=.1cm of hw](hos){ホストOS};
            \node[str,fill=gray!40,text width=.9\textwidth,above=.1cm of hos](hsw){ホスト型仮想化ソフトウェア};
            \node[str,fill=gray!30,text width=.43\textwidth,above=.1cm of hsw.north west,anchor=south west](gos1){ゲストOS};
            \node[str,fill=gray!30,text width=.43\textwidth,above=.1cm of hsw.north east,anchor=south east](gos2){ゲストOS};
            \node[str,fill=gray!20,text width=.2\textwidth,above=.1cm of gos1.north west,anchor=south west](sw1){\tiny アプリケーション};
            \node[str,fill=gray!20,text width=.2\textwidth,above=.1cm of gos1.north east,anchor=south east](sw2){\tiny アプリケーション};
            \node[str,fill=gray!20,text width=.2\textwidth,above=.1cm of gos2.north west,anchor=south west](sw3){\tiny アプリケーション};
            \node[str,fill=gray!20,text width=.2\textwidth,above=.1cm of gos2.north east,anchor=south east](sw4){\tiny アプリケーション};
        \end{tikzpicture}
    \end{figure}
    % \hyperlink{ハイパーバイザ型}{\beamergotobutton{ハイパーバイザ型}}\hypertarget{ホスト型}{}
\end{frame}
\begin{frame}[t]{\ftitle}
    \begin{exampleblock}{ホスト型仮想化ソフトウェア 例}
        \begin{minipage}[b]{.8\textwidth}
            \begin{itemize}
                \setlength{\itemsep}{1em}
                \item VMware Workstation Player
                \item VMware Fusion
                \item Oracle VM Virtualbox
            \end{itemize}
        \end{minipage}
        \begin{minipage}[b]{.15\textwidth}
            \centering
            \includegraphics[keepaspectratio,width=\textwidth]{virtualbox_logo.png}\\
            {\tiny Virtualbox\cite{VMBox}}
        \end{minipage}
    \end{exampleblock}
    \begin{minipage}[t]{.48\textwidth}
        \textbf{メリット}
        \begin{itemize}
            \item 既存マシンが利用できる点．
            \item 仮想化に必要なソフトウェアが扱いやすい．
        \end{itemize}
    \end{minipage}
    \begin{minipage}[t]{.48\textwidth}
        \textbf{デメリット}
        \begin{itemize}
            \item ホストOSを動作させるための物理リソースが必要．
        \end{itemize}
    \end{minipage}\\
    \hfill\cite{itmanage}
\end{frame}
\subsection{ハイパーバイザ型}
\begin{frame}[t]{\ftitle}
    ハイパーバイザとは「仮想化のためのOS」のようなもの．
    \begin{figure}[b]
        \centering
        \begin{tikzpicture}
            \node[str,fill=gray!60,text width=.9\textwidth](hw){ハードウェア};
            \node[str,fill=gray!45,text width=.9\textwidth,above=.1cm of hw,minimum height=2.1cm](hpv){ハイパーバイザ};
            \node[str,fill=gray!30,text width=.43\textwidth,above=.1cm of hpv.north west,anchor=south west](gos1){ゲストOS};
            \node[str,fill=gray!30,text width=.43\textwidth,above=.1cm of hpv.north east,anchor=south east](gos2){ゲストOS};
            \node[str,fill=gray!20,text width=.2\textwidth,above=.1cm of gos1.north west,anchor=south west](sw1){\tiny アプリケーション};
            \node[str,fill=gray!20,text width=.2\textwidth,above=.1cm of gos1.north east,anchor=south east](sw2){\tiny アプリケーション};
            \node[str,fill=gray!20,text width=.2\textwidth,above=.1cm of gos2.north west,anchor=south west](sw3){\tiny アプリケーション};
            \node[str,fill=gray!20,text width=.2\textwidth,above=.1cm of gos2.north east,anchor=south east](sw4){\tiny アプリケーション};
        \end{tikzpicture}
    \end{figure}
    % \hyperlink{ホスト型}{\beamergotobutton{ホスト型}}\hypertarget{ハイパーバイザ型}{}
\end{frame}
\begin{frame}[t]{\ftitle}
    \begin{exampleblock}{ハイパーバイザ 例}
        \begin{itemize}
            \item VMware ESXi
            \item Linux KVM
            \item Microsoft Hyper-V
        \end{itemize}
    \end{exampleblock}
    \begin{minipage}[t]{.48\textwidth}
        \textbf{メリット}
        \begin{itemize}
            \item システム全体の観点から見てリソースの使用効率がよい．
            \item ホストOSが不要でハードウェアを直接制御が可能．
        \end{itemize}
    \end{minipage}
    \begin{minipage}[t]{.48\textwidth}
        \textbf{デメリット}
        \begin{itemize}
            \item 仮想化環境の高度な管理を実現するツールが標準装備されていない場合がある．
        \end{itemize}
    \end{minipage}\\
    \hfill\cite{itmanage}
\end{frame}
\subsection{コンテナ}
\begin{frame}[t]{\ftitle}
    ``アプリケーションを実行するための領域（ユーザ空間）を複数に分割して利用するもの''\cite{itmanage}．
    \begin{figure}[b]
        \centering
        \begin{tikzpicture}
            \node[str,fill=gray!60,text width=.9\textwidth](hw){ハードウェア};
            \node[str,fill=gray!50,text width=.9\textwidth,above=.1cm of hw](hos){ホストOS};
            \node[str,fill=gray!40,text width=.9\textwidth,above=.1cm of hos](hsw){Docker または Kubernetes};
            \node[str,fill=gray!30,text width=.2\textwidth,above=.1cm of hsw.north west,anchor=south west](gos1){ミドルウェア};
            \node[str,fill=gray!30,text width=.2\textwidth,above=.1cm of hsw.north east,anchor=south east](gos2){ミドルウェア};
            \node[str,fill=gray!20,text width=.2\textwidth,above=.1cm of gos1.north west,anchor=south west](sw1){\tiny アプリケーション};
            \node[str,fill=gray!20,text width=.2\textwidth,above=.1cm of gos2.north east,anchor=south east](sw4){\tiny アプリケーション};
            \node[inner sep=.5mm,fit={(gos1)(sw1)},draw,thick,dotted,rounded corners](wrap1){};
            \node[inner sep=.5mm,fit={(gos2)(sw4)},draw,thick,dotted,rounded corners](wrap2){};
            \node at ($(wrap1)!0.5!(wrap2)+(0,.5cm)$)(cont){コンテナ};
            \draw[-latex](cont.west)--(wrap1.east);
            \draw[-latex](cont.east)--(wrap2.west);
        \end{tikzpicture}
    \end{figure}
\end{frame}
\section{Docker（Option）}\label{chap:docker}
\tocc
\subsection{コンテナ立ち上げと利用}
\begin{frame}[t]{\ftitle}
    \begin{itemize}
        \setlength{\itemsep}{1em}
        \item コンテナを利用する．
              \begin{enumerate}
                  \setlength{\itemsep}{.5em}
                  \item イメージを作成（\texttt{build}）または，取得（\texttt{docker pull}）する．\\
                        イメージの作成には\texttt{Dockerfile} を書く．
                  \item コンテナを作成して起動（\texttt{run}）する．
                  \item コンテナ内でコマンドを実行（\texttt{exec}）する．
              \end{enumerate}
        \item コンテナ利用後．
              \begin{enumerate}
                  \setlength{\itemsep}{.5em}
                  \item コンテナを停止（\texttt{stop}）する．
                  \item コンテナを削除（\texttt{rm}）する．
                  \item イメージを削除（\texttt{rmi}）する．
              \end{enumerate}
    \end{itemize}
\end{frame}
\subsection{イメージ作成・取得}
\begin{frame}[containsverbatim,t]{\ftitle}
    \begin{itemize}
        \item \texttt{Dockerfile}
              \begin{lstlisting}[language=docker-compose-2]
FROM mysql:latest
            \end{lstlisting}
              \begin{lstlisting}[backgroundcolor={\color{black}},basicstyle={\color{white}\ttfamily},]
$ docker build ./ -t mysql
                            \end{lstlisting}
        \item レジストリからイメージを取得（Pull）する．
              \begin{lstlisting}[backgroundcolor={\color{black}},basicstyle={\color{white}\ttfamily},]
$ docker pull mysql:latest
            \end{lstlisting}
    \end{itemize}
\end{frame}
\subsection{実演}
\begin{frame}{\ftitle}
    \includegraphics[keepaspectratio,width=\textwidth]{docker_logo.png}\\\vspace{1em}
    \hfill\cite{docker}
\end{frame}
\section{準仮想化と完全仮想化}\hypertarget{sec:準仮想化と完全仮想化}{}
\tocc
\begin{frame}[t]{\fftitle}
    ホストOSとハイパーバイザ型を比較すると，性能的にはハイパーバイザ型が有利．
    ハイパーバイザの中にも\textbf{完全仮想化}と\textbf{準仮想化}がある．
    \begin{block}{完全仮想化}
        ハードウェアも含めて，すべてを仮想化する方式．ハードウェアも論理ハードウェア．
    \end{block}
    \begin{block}{準仮想化}
        ゲストOSを少し改造して，ゲストOSから直接ハードウェアを操作できる方式．
        処理速度が大きくなる．（理由は次頁．）
    \end{block}
    \hfill\hyperlink{sec:クラウドコンピューティング}{\beamerreturnbutton{詳しい説明を飛ばす}}
\end{frame}
\begin{frame}[t]{\fftitle （詳しく）}
    \begin{block}{完全仮想化}
        ハイパーバイザ上で動作するゲストOSの中身に何も変更を加えないで動作させる方式．
        ゲストOSからハードウェアを操作する（特権命令の実行）は許されないので，ハイパーバイザで処理を変換する．
    \end{block}
    \begin{block}{準仮想化}
        ゲストOSに変更を加えて性能向上を図った方式．
        ゲストOSからハードウェアを操作する部分（特権命令）を，ハイパーバイザに直接処理が渡るように変更を加える．
        ハイパーバイザでの命令変換時にかかるオーバーヘッドを削減する
    \end{block}
    \hfill\cite{仮想化技術}
\end{frame}
\begin{frame}[t]{\fftitle （詳しく）}
    \newcommand{\karano}{からの}
    命令\texttt{A}：ゲストOSからの特権命令\footnote{CPU のステータスを変更するといった，特別な権限を与えられたプログラムしか使用できない命令．}．\\
    命令\texttt{X}：ハイパーバイザからの特権命令．
    \begin{figure}
        \centering
        \begin{tikzpicture}
            \node[draw,fill=gray!60,minimum height=1cm,align=center,text width=.4\textwidth](vm1){命令\texttt{A}};
            \node[draw,fill=gray!40,minimum height=1cm,align=center,text width=.4\textwidth,below=.5cm of vm1](vmm1){命令\texttt{A}\ \raisebox{.25\baselineskip}{\tikz{\draw[-latex](0,0)--(.5,0)}}\ \texttt{X}};
            \node[draw,fill=gray!20,minimum height=1cm,align=center,text width=.4\textwidth,below=.5cm of vmm1](hw1){命令\texttt{X}};
            \node[anchor=north west,draw,fill=white] at (vm1.north west){\tiny ゲストOS};
            \node[anchor=north west,draw,fill=white] at (vmm1.north west){\fontsize{4pt}{0pt}\selectfont ハイパーバイザ};
            \node[anchor=north west,draw,fill=white] at (hw1.north west){\tiny ハードウェア};
            \node[draw,fill=gray!60,minimum height=1cm,align=center,text width=.4\textwidth,right=.5cm of vm1](vm2){命令\texttt{X}};
            \node[draw,fill=gray!40,minimum height=1cm,align=center,text width=.4\textwidth,below=.5cm of vm2](vmm2){命令\texttt{X}};
            \node[draw,fill=gray!20,minimum height=1cm,align=center,text width=.4\textwidth,below=.5cm of vmm2](hw2){命令\texttt{X}};
            \node[anchor=north west,draw,fill=white] at (vm2.north west){\tiny ゲストOS};
            \node[anchor=north west,draw,fill=white] at (vmm2.north west){\fontsize{4pt}{0pt}\selectfont ハイパーバイザ};
            \node[anchor=north west,draw,fill=white] at (hw2.north west){\tiny ハードウェア};
            \node[below] at (hw1.south){完全仮想化};
            \node[below] at (hw2.south){準仮想化};
            \foreach \u \v in {vm1/vmm1,vmm1/hw1,vm2/vmm2,vmm2/hw2}{
                    \draw[-latex](\u)--(\v);
                }
            \node[align=center,draw,fill=gray!60,rounded corners] at($(vm2.west)+(.6cm,-.5cm)$) {\tiny あらかじめ特権命\\\tiny 令を変換しておく};
            \node[align=center,draw,fill=gray!40,rounded corners] at($(vmm1.east)+(-.8cm,-.5cm)$) {\tiny 特権命令を変換する};
        \end{tikzpicture}
    \end{figure}
\end{frame}
\section{クラウドコンピューティング}\hypertarget{sec:クラウドコンピューティング}{}
\tocc
\subsection{クラウドとは}
\begin{frame}[t]{\ftitle}
    \begin{block}{クラウドコンピューティング}
        クラウドコンピューティング（以下，クラウド）とは，コンピュータリソース尾の利用形態．
        コンピュータの計算リソースやストレージ領域，アプリケーションによる処理をネットワーク経由で提供する．\hfill\cite{2015amazon}
    \end{block}
    クラウドは以下の4つに分類されることが多い．
    \begin{itemize}
        \setlength{\itemsep}{1em}
        \item IaaS（Infrastructure as a Service）
        \item PaaS（Platform as a Service）
        \item SaaS（Software as a Service）
        \item DaaS（Desktop as a Service）% VDIとの違い
    \end{itemize}
\end{frame}
\subsection{クラウドの分類}
\begin{frame}[t]{\ftitle}
    \begin{itemize}
        \item \textbf{IaaS}\\
              仮想サーバやストレージ，ネットワークサービスなどをインターネット経由で提供する．\\（AWS，Azure，GCP\footnote{Google Cloud Pratform}，IBM Cloudなど）
        \item \textbf{PasS}\\
              ミドルウェアをサービスとして提供する．
              OSとミドルウェアの管理はサービス提供者で行われ，ユーザはミドルウェアを直接操作できる．\\（AWS，Azure，GCPなど）
        \item \textbf{SaaS}\\
              ソフトウェアをサービスとして提供する．\\（Slack，Zoom，Google Calendarなど）
    \end{itemize}
    \hfill\cite{2015amazon}
\end{frame}
\begin{frame}[t]{\ftitle}
    \begin{itemize}
        \item \textbf{DaaS}\\
              ``クラウドサービスとして提供される仮想デスクトップ環境を組織として利用するしかけ．''\hfill\cite{ISディジタル辞典DaaS}\\
              %   \begin{itemize}
              %       \item[{\bfseries\Large\cmark}] \underline{クライアントを安価に抑えることができる．}（クライアントが依頼した処理をサーバが処理するため）
              %       \item[{\bfseries\Large\xmark}] \underline{ネットワークネットワーク負荷が大きくなる．}（ディスクトップ環境を提供するため）
              %   \end{itemize}
              （Azure Virtual Desktopなど）\\
              \hrulefill
        \item \textbf{VDI}（Virtual Desktop Infrastructure）\\
              ``クライアントPCのディスクトップサーバ上に仮想化して集約し，稼働させる仕組み．''\hfill\cite{ISディジタル辞典VDI}
              \item[\textbf{\large !}]DaaSはVDIの一種である．
    \end{itemize}
\end{frame}
\subsection{代表的なIaaS}
\begin{frame}[t]{\ftitle}
    \begin{itemize}
        \setlength{\itemsep}{1em}
        \item \underline{AWS}
        \item Microsoft Azure
        \item Google Cloud Pratform
        \item さくらのクラウド（さくらインターネット）
    \end{itemize}
\end{frame}
\section{AWS}
\tocc
\subsection{AWSとは}
\begin{frame}[t]{\ftitle}
    \begin{block}{AWS}
        Amazon Web Service の略称．Amazonが提供するクラウドサービスで，ネットワークを経由して仮想コンピュータやストレージなどのサービスを提供している．\hfill\cite{2015amazon}
    \end{block}
    \begin{minipage}{\textwidth}
        \centering
        \begin{minipage}[t]{.32\textwidth}
            \centering
            \includegraphics[keepaspectratio,width=\textwidth]{Arch_Amazon-EC2_64@5x.png}\\
            {\tiny Amazon Elastic Compute Cloud (Amazon EC2)\cite{aws}}
        \end{minipage}
        \begin{minipage}[t]{.32\textwidth}
            \centering
            \includegraphics[keepaspectratio,width=\textwidth]{Arch_Amazon-Simple-Storage-Service_64@5x.png}
            {\tiny Amazon Simple Storage Service (Amazon S3)\cite{aws}}
        \end{minipage}
        \begin{minipage}[t]{.32\textwidth}
            \centering
            \includegraphics[keepaspectratio,width=\textwidth]{Arch_Amazon-RDS_64@5x.png}\\
            {\tiny Amazon Relational Database Service (Amazon RDS)\cite{aws}}
        \end{minipage}
    \end{minipage}
\end{frame}
\subsection{オンプレミスとAWS}
\begin{frame}[t]{\ftitle}
    リソースを効率よく運用できる．\hfill\cite{2015amazon}
    \begin{figure}
        \centering
        \begin{minipage}[b]{.48\textwidth}
            \centering
            \begin{tikzpicture}
                \coordinate (O) at (0,0);
                \coordinate (A) at (1,2);
                \coordinate (B) at (2,1.5);
                \coordinate (C) at (3,3);
                \coordinate (D) at (4,2);
                \draw[red,thick, name path=curve, save path=\curve]plot[smooth, tension=0.7]coordinates {(O)(A)(B)(C)(D)};
                \draw[name path=line,thick,dashed](0,2.5)--(4,2.5);
                \draw[thick,dashed](0,2.5)--(4,2.5);
                \draw[thick,-latex](0,0)--(4,0);
                \draw[thick,-latex](0,0)--(0,4.5);
                \node[below] at ($(0,0)!0.5!(4,0)$){\scriptsize 時間};
                \node[left] at ($(0,0)!0.5!(0,4.5)$){\scriptsize\rotatebox{90}{リソース量}};
                \node[below,text=red]at(B){\scriptsize 実際のリソース};
                \node[above]at(C){\scriptsize 機会損失};
                \node[below right]at(0,2.5){\scriptsize 無駄になる};
            \end{tikzpicture}\\
            \vspace{1em}
            オンプレミス
        \end{minipage}
        \begin{minipage}[b]{.48\textwidth}
            \centering
            \begin{tikzpicture}
                \coordinate (O) at (0,0);
                \coordinate (A) at (1,2);
                \coordinate (B) at (2,1.5);
                \coordinate (C) at (3,3);
                \coordinate (D) at (4,2);
                \newcommand{\avv}{0.2cm}
                \draw[red,thick]plot[smooth, tension=0.7]coordinates {(O)(A)(B)(C)(D)};
                \draw[thick,dashed]plot[smooth,tension=0.7] coordinates {($(O)+(0,\avv)$)($(A)+(0,\avv)$)($(B)+(0,\avv)$)($(C)+(0,\avv)$)($(D)+(0,\avv)$)};
                \draw[thick,-latex](0,0)--(4,0);
                \draw[thick,-latex](0,0)--(0,4.5);
                \node[below] at ($(0,0)!0.5!(4,0)$){\scriptsize 時間};
                \node[left] at ($(0,0)!0.5!(0,4.5)$){\scriptsize\rotatebox{90}{リソース量}};
                \node[above]at($(C)+(0,\avv)$){\scriptsize 用意したリソース};
                \node[below,text=red]at(B){\scriptsize 実際のリソース};
            \end{tikzpicture}\\
            \vspace{1em}
            AWSの利用
        \end{minipage}
    \end{figure}
\end{frame}
\subsection{レンタルサーバとAWS}
\begin{frame}[t]{\ftitle}
    \begin{block}{レンタルサーバ}
        ``1台のサーバを複数のユーザが共用で利用する形態''\hfill\cite{2015amazon}
    \end{block}
    \begin{table}
        \fontsize{10pt}{10pt}\selectfont
        \centering
        \renewcommand{\arraystretch}{1.3}
        \newcommand{\hobo}{ほぼ}
        \begin{tabularx}{\textwidth}{cRRR}
                             & \multicolumn{1}{c}{共用サーバ} & \multicolumn{1}{c}{専用サーバ} & \multicolumn{1}{c}{仮想専用サーバ}          \\
            \hline
            {\tiny 利用形態}     & 1台の物理サーバを分割して利用           & 1台の物理サーバを占有               & 1台の物理サーバ上に立てた仮想サーバを利用                \\
            \hline
            {\tiny コスト}      & 安い                        & 高い                        & \textcolor{red}{中間}                  \\
            \hline
            {\tiny カスタマイズ}   & \hobo 不可能                 & 可能                        & \textcolor{red}{可能}                  \\
            \hline
            {\tiny セキュリティ}   & コントロール\textbf{不可}         & 高くコントロール\textbf{可能}       & \textcolor{red}{高くコントロール\textbf{可能}} \\
            \hline
            {\tiny 他のユーザの影響} & 強く受ける                     & 受けない                      & 多少受ける                                \\
            \hline
        \end{tabularx}
    \end{table}
    \hfill 引用\cite{2015amazon}
\end{frame}
\subsection{AWS導入例}
\begin{frame}[t]{\ftitle}
    \begin{itemize}
        \item Cloud LaTeX（アカリク）
              \begin{itemize}
                  \item コンパイルなど（EC2）
                  \item プロジェクト，テンプレート管理（S3，MySQL）
              \end{itemize}
        \item 株式会社任天堂
              \begin{itemize}
                  \item 『マリオカート ツアー』のDBに，「Amazon Aurora」（RDMS） を採用．
              \end{itemize}
        \item 株式会社東京海上日動火災保険
        \item JCB
        \item 盛岡市
        \item 浜松市
        \item SUBARU
    \end{itemize}
    \hfill\cite{導入事例}
\end{frame}
\begin{frame}[allowframebreaks]{Reference}
    \bibliography{bib}
\end{frame}
\end{document}